Since both sides of the relation to be proved are additive, we may suppose that @@M0@@ is the class of an @@M1@@-Module locally free of finite type @@M2@@. The element @@M3@@ is then the class of the complex @@M4@@, obtained by taking a projective resolution of the complex @@M5@@ and tensoring it with @@M6@@. The homology sheaves of this complex are therefore the @@M7@@. By 2.4, we have the isomorphism:
@@M8@@
By 2.5, we have the isomorphism:
@@M9@@

The @@M10@@ are therefore @@M11@@-Modules locally free of finite type, and the class in @@M12@@ of the complex @@M13@@ is the alternating sum of the classes of its homology sheaves, namely:
@@M14@@

\sgaStableTarget{sga6.e07.cor.2-8}\textbf{Corollary 2.8.} \uline{Under the hypotheses of} 2.7, \uline{let} @@M15@@; \uline{we have the relation}:
\sgaStableTarget{sga6.e07.eq.2-8-1}@@M16@@

By 2.7, @@M17@@.

Applying the projection formula, we therefore find:
@@M18@@

\sgaStableTarget{sga6.e07.sec.3}\section*{3. Computation of \texorpdfstring{@@M19@@}{K} of a blown-up scheme}

\sgaStableTarget{sga6.e07.para.3-1}\textbf{3.1.} In this section, @@M20@@ denotes a regular closed immersion. Denote by @@M21@@ the scheme obtained by blowing up @@M22@@, and put @@M23@@.
